% Options for packages loaded elsewhere
\PassOptionsToPackage{unicode}{hyperref}
\PassOptionsToPackage{hyphens}{url}
%
\documentclass[
]{article}
\usepackage{amsmath,amssymb}
\usepackage{lmodern}
\usepackage{iftex}
\ifPDFTeX
  \usepackage[T1]{fontenc}
  \usepackage[utf8]{inputenc}
  \usepackage{textcomp} % provide euro and other symbols
\else % if luatex or xetex
  \usepackage{unicode-math}
  \defaultfontfeatures{Scale=MatchLowercase}
  \defaultfontfeatures[\rmfamily]{Ligatures=TeX,Scale=1}
\fi
% Use upquote if available, for straight quotes in verbatim environments
\IfFileExists{upquote.sty}{\usepackage{upquote}}{}
\IfFileExists{microtype.sty}{% use microtype if available
  \usepackage[]{microtype}
  \UseMicrotypeSet[protrusion]{basicmath} % disable protrusion for tt fonts
}{}
\makeatletter
\@ifundefined{KOMAClassName}{% if non-KOMA class
  \IfFileExists{parskip.sty}{%
    \usepackage{parskip}
  }{% else
    \setlength{\parindent}{0pt}
    \setlength{\parskip}{6pt plus 2pt minus 1pt}}
}{% if KOMA class
  \KOMAoptions{parskip=half}}
\makeatother
\usepackage{xcolor}
\setlength{\emergencystretch}{3em} % prevent overfull lines
\providecommand{\tightlist}{%
  \setlength{\itemsep}{0pt}\setlength{\parskip}{0pt}}
\setcounter{secnumdepth}{-\maxdimen} % remove section numbering
\ifLuaTeX
  \usepackage{selnolig}  % disable illegal ligatures
\fi
\IfFileExists{bookmark.sty}{\usepackage{bookmark}}{\usepackage{hyperref}}
\IfFileExists{xurl.sty}{\usepackage{xurl}}{} % add URL line breaks if available
\urlstyle{same} % disable monospaced font for URLs
\hypersetup{
  hidelinks,
  pdfcreator={LaTeX via pandoc}}

\author{Dietmar Gerald Schrausser}
\date{09.05.2026}


\begin{document}

\hypertarget{foliumblat-vv}{%
\section{Foliu{[}m{]}/Blat Vv}\label{foliumblat-vv}}

Eine zeitgenössische deutsche Übersetzung wird präsentiert, die auf
einem Vergleich der bestehenden lateinischen, frühneuhochdeutschen
(ENHG) und englischen (Übersetzungs-) Texte basiert, um ein besseres
Verständnis zu ermöglichen. Insbesondere da der ENHG-Text im Vergleich
zum modernen Hochdeutschen schwer verständlich ist, da dieser fundierte
Deutschkenntnisse (als Muttersprache) sowie Kenntnisse verschiedener
Dialekte voraussetzt.

\hypertarget{zur-heiligung-des-siebenten-tages}{%
\subsection{Zur Heiligung des siebenten
Tages}\label{zur-heiligung-des-siebenten-tages}}

Als die Welt, ein Werk \emph{göttlicher Weisheit}, in sechs Tagen
vollendet war, Himmel und Erde erschaffen, geordnet und schliesslich
geschmückt, da beendete der glorreiche Gott sein Werk und ruhte am
\emph{siebten} Tag von seinen (Hand-)Werken. Diese Komplettierung,
bedeutete jedoch nicht ein \emph{Beenden}, als wäre er der Arbeit
\emph{müde}, sondern war vielmehr um \emph{kein}\footnote{`facere
  cessauit' und `zemachen {[}\ldots{]} nit vergange{[}n{]}'.} weiteres
Geschöpf anderer Materie oder Ebenbilder zu schaffen, zumal das Werk der
Schöpfung nicht zu wirken aufhört.\footnote{aus Petrus Lombardus'
  `Sententiarum' (Petrus,
  \href{https://doi.org/10.3931/e-rara-18241}{1484}, Buch II, Kapitel
  7).}

Der Herr segnete und heiligte \emph{diesen} Tag und nannte ihn
\emph{Sabbatum}, was auf Hebräisch \emph{Ruhe}\footnote{`requiem' und
  `ein rüe'.} bedeutet, so wie er an diesem Tag von all seinem Werk
ruhte.

\begin{quote}
Eo q{[}uia{]} in ipso cessauerat ab om{[}n{]}i opere q{[}uo{]}d
patrarat.\footnote{`Denn an diesem hatte er mit all seinen Werken
  aufgehört', `Vulgata', eine Variation von Genesis 2,3 (vgl. Jiménez de
  Cisneros, \href{https://doi.org/10.3931/e-rara-46695}{1517}).}
\end{quote}

Es ist auch unter den Juden üblich diesen Tag zu feiern indem sie ihrer
Arbeit entsagen, ebenso war es \emph{Gesetz} vieler heidnischer Völker
diesen Tag zu feiern.

Damit sind wir nun am Ende der \emph{göttlichen Werke} angelangt.
Entsprechend sollten wir (i) dem \emph{mit liebevoller Ehrfurcht
begegnen},\footnote{vielleicht besser als `fürchten, lieben und ehren'.}
in dem alles Sichtbare und Unsichtbare existiert und vom Herrn des
Himmels, dem Herrn aller guten Dinge, dem die Macht im Himmel und auf
Erden gegeben, sowohl (ii) gegenwärtiges \emph{Glück}\footnote{`presentia
  bona: quatenus bona sint' und `die gegenwürtigen güter. sover die gut
  sind', bezieht sich auf \emph{gegenwärtige Segnungen}, in einem
  philosophischen bzw. theologischen Kontext, gegenwärtig, materielle
  oder irdische Dinge, im Gegensatz zu ewigen, geistigen Gütern (vgl.
  Thomas von Aquin, \href{https://doi.org/10.3931/e-rara-15170}{1485}
  oder Augustinus,
  \href{https://archive.org/details/expositionsonboo39augu/mode/1up}{1847}).}
(soweit es \emph{gut}\footnote{\emph{wohlwollend}.} ist) als auch die
wahre \emph{Seligkeit} des ewigen Lebens \emph{ersehnen}\footnote{`queramus'
  und `suchen'.}.

\hypertarget{references}{%
\subsection{References}\label{references}}

Augustinus, A. (1847). \emph{Expositions on the Book of Psalms}. Oxford:
J. H. Parker.
\url{https://archive.org/details/expositionsonboo39augu/mode/1up}

Jiménez de Cisneros, F. (1517). \emph{Biblia Polyglotta Complutensis}.
Complutum: Arnaldo Guillén de Brocar.
\url{https://doi.org/10.3931/e-rara-46695}

Petrus. (1484). \emph{Sententiarum Libri IV}. Basel.
\url{https://doi.org/10.3931/e-rara-18241}

Thomas. (1485). \emph{Summa Theologica: Partes 1-3}. Basel: Michael
Wenssler. \url{https://doi.org/10.3931/e-rara-15170}

\end{document}
